%--------------------------
\begin{frame}{Question 2}

\begin{block}{}
Les revendications :

\begin{itemize}
  \item[$\bullet$] a. peuvent faire l’objet de modifications après réception du rapport de recherche 
  internationale.  
  \item[$\bullet$] b. ne sont pas la seule partie de la demande internationale qui peut faire l’objet de 
  modifications pendant la phase internationale. 
  \item[$\circ$] c. peuvent être une partie manquante de la demande internationale lors du dépôt sans 
  que cela empêche nécessairement l’attribution d’une date de dépôt international. 
  \item[$\circ$] d. Aucune des réponses a, b ou c n’est correcte. 
\end{itemize}

\end{block}

\end{frame}

\begin{frame}{Question 2}

\textbf{Base juridique :} \pctart{19} + \pctart{34} + \pctart{11}(1).  

\begin{itemize}
  \item[$\bullet$] \textbf{a.} \; Correct : les revendications peuvent être modifiées après réception de l’ISR,
  dans le délai prévu (\pctart{19} + \pctrule{46}).

  \item[$\bullet$] \textbf{b.} \; Correct : la description et les dessins peuvent aussi être modifiés dans le cadre
  du chapitre II (\pctart{34}).

  \item[$\bullet$] \textbf{c.} \; Faux : au moins une revendication
est requise pour obtenir une date de dépôt international
(\pctart{11}.1.iii.e)).

  \item[$\circ$] \textbf{d.} \; Faux puisque \textbf{a}, \textbf{b} et \textbf{c} sont correctes.
\end{itemize}

\end{frame}
